% Sample deck for the presenter console.
% Double-width pages (slide | notes) + an inline animate animation driven
% by pgfplots, so the full pipeline (halves rendering, CropBox-derived
% audience deck, animations in the official pdf.js viewer) can be
% demonstrated without any external files.  Build: pdflatex sample.tex
\documentclass[aspectratio=169]{beamer}
\usepackage{animate}
\usepackage{pgfplots}
\pgfplotsset{compat=1.18}

\setbeameroption{show notes on second screen=right}
% Bare note pages: no beamer chrome, small font (same layout idea as the
% real deck this console was built for).
\setbeamertemplate{note page}{%
\vfill\footnotesize\raggedright\insertnote\par\vfill}

\begin{document}

\begin{frame}{Presenter console --- sample deck}
  This PDF is \emph{double-width}: the slide occupies the left half of
  every page, the speaker notes the right half.
  \note{Welcome! This is the notes half --- only you see it. The
    audience window receives a cropped slide-only copy of this same
    PDF. Turn pages with the arrow keys or the buttons; the timer
    starts on your first page turn (reset it with the button next to
    it). Drag the separators to resize the three panes; double-click a
  separator to reset.}
\end{frame}

\begin{frame}{Animated pgfplots: policy training}
  \begin{center}
    % 36 frames @ 12fps = 3 s loop.  The animate package has NO built-in
    % frame transitions (hard cuts only), so fluid fades are baked into
    % the frames themselves:
    %   - data fades in over frames 0-5 (appearing transition),
    %   - holds at full opacity,
    %   - fades out over frames 32-35 so the loop wrap blends into the
    %     fade-in of frame 0 (seamless loop, no hard cut),
    %   - a dot rides each curve frontier to accent the motion.
    % poster=31: the static appearance (console preview pane, printing,
    % viewers without JS) shows a nearly complete chart instead of the
    % blank first fade frame.
    \begin{animateinline}[autoplay,loop,poster=31]{12}
      \multiframe{36}{i=0+1}{%
        \begin{tikzpicture}
          \pgfmathsetmacro{\af}{\i<6 ? \i/5 : (\i>31 ? (35-\i)/3 : 1)}
          \pgfmathsetmacro{\X}{100*(\i+1)/36}
          \pgfmathsetmacro{\yA}{1 - 2*exp(-\X/38) + 0.1*sin(deg(3*\X/10))}
          \pgfmathsetmacro{\yB}{-1 + 2*(1 - exp(-\X/55)) + 0.08*cos(deg(\X/12))}
          \begin{axis}[
              width=0.72\textwidth,
              height=0.62\textheight,
              xlabel={environment step},
              ylabel={mean reward},
              xmin=0, xmax=100,
              ymin=-1.2, ymax=1.2,
              grid=both,
              major grid style={gray!30},
              legend pos=south east,
              legend style={font=\scriptsize},
            ]
            % Smooth learning curves, progressively revealed with a
            % fade-in/out envelope (\af) for fluid appearing transitions.
            \addplot[thick, blue, domain=0:\X, samples=50,
            smooth, opacity=\af]
            {1 - 2*exp(-x/38) + 0.1*sin(deg(3*x/10))};
            \addlegendentry{policy}
            \addplot[dashed, orange, domain=0:\X, samples=50,
            smooth, opacity=\af]
            {-1 + 2*(1 - exp(-x/55)) + 0.08*cos(deg(x/12))};
            \addlegendentry{baseline}
            % Frontier dots riding the leading edge of each curve.
            \addplot[only marks, mark=*, mark size=1.5pt,
            blue, opacity=\af] coordinates {(\X,\yA)};
            \addplot[only marks, mark=*, mark size=1.5pt,
            orange, opacity=\af] coordinates {(\X,\yB)};
          \end{axis}
      \end{tikzpicture}}
    \end{animateinline}
  \end{center}
  \note{This is a pgfplots chart animated with the animate package: the
    two training curves are progressively revealed, like watching a run
    progress. The motion is kept fluid by baking opacity ramps into the
    frames (the animate package has no built-in frame transitions): the
    data fades in at the start of each loop, fades out at the end ---
    so the loop wraps without a hard cut --- and a dot rides each
    curve's frontier. The audience deck is derived with pdf-lib CropBox
    surgery, which keeps the widget annotations --- so the animation
  still plays in the audience window (official pdf.js viewer).}
\end{frame}

\begin{frame}{Pick your own deck}
  Use \textbf{Browse PDF\dots} to open your own double-width build ---
  for example beamer's \texttt{show notes on second screen=right}.
  \note{Slides without a \texttt{\textbackslash note\{\}} show an empty
    right half in the console --- that is content, not a bug. The
  next-slide pane on the far right previews what comes next.}
\end{frame}

\end{document}
